\documentclass[11pt]{article}
\usepackage[margin=2.0cm]{geometry}
\usepackage{booktabs}
\usepackage{amsmath}
\usepackage{amssymb}
\usepackage{enumitem}
\usepackage{array}
\usepackage{listings}
\usepackage{xcolor}
\usepackage[T1]{fontenc}
\usepackage{lmodern}

\setlength{\parindent}{0pt}
\setlength{\parskip}{3.6pt}
\newcommand{\code}[1]{\texttt{#1}}
\newcommand{\hr}{\vspace{2pt}\hrule\vspace{6pt}}
\renewcommand{\arraystretch}{1.12}
\newcommand{\ok}{\checkmark}
\newcommand{\bad}{\textbf{FAIL}}

\definecolor{kw}{RGB}{0,90,160}
\definecolor{cmt}{RGB}{110,110,110}
\lstdefinestyle{vhdl}{
  language=VHDL, basicstyle=\ttfamily\footnotesize,
  keywordstyle=\color{kw}\bfseries, commentstyle=\color{cmt}\itshape,
  showstringspaces=false, columns=fullflexible, breaklines=true,
  frame=leftline, framesep=6pt, xleftmargin=8pt, aboveskip=6pt, belowskip=6pt
}
\lstset{style=vhdl}

\begin{document}

\begin{center}
{\Large\bfseries EE\,362 --- Computer Organization Laboratory}\\[2pt]
{\large\bfseries Lab 1 --- Guided Examples: Solution}\\[2pt]
{\itshape Example 1 Majority-of-3 \quad$\cdot$\quad Example 2 Transparent Latch \quad$\cdot$\quad Example 3 Sensitivity List}\\[4pt]
Fall 2026 \quad$\cdot$\quad Lab TA: Gandhi Torres
\end{center}
\hr

Three examples, and each one has exactly one job:

\begin{center}
\begin{tabular}{clll}
\toprule
& circuit & new idea & feeds the assignment as \\
\midrule
1 & \code{majority3} & the \textbf{tool flow} and the \textbf{testbench} & the check-and-report pattern \\
2 & \code{trans\_latch} & \textbf{memory}: $Q$ on the right-hand side & \code{TransLatch} minus \code{Clr} \\
3 & \code{trans\_latch\_bad} & \textbf{one missing name} & the bug planted in your starter \\
\bottomrule
\end{tabular}
\end{center}

Every result quoted below --- transcripts, warnings, register counts, timing --- was
measured with Quartus Prime Lite 25.1std and Questa on device \code{5CSEMA4U23C6}.
Nothing is paraphrased from memory.

\vspace{4pt}\hr

\section*{0. What to enter in Quartus --- once, for all three}

Each example folder has its own \code{.qpf}. Open it (\code{File $\rightarrow$ Open
Project}) and confirm these settings before compiling --- a wrong top-level entity
is the most common reason ``nothing happens''.

\begin{center}
\small
\begin{tabular}{>{\raggedright\arraybackslash}p{2.4cm}>{\raggedright\arraybackslash}p{7.2cm}>{\raggedright\arraybackslash}p{5.6cm}}
\toprule
setting & where & value \\
\midrule
Family / device & \code{Assignments $\rightarrow$ Device} & Cyclone V, \code{5CSEMA4U23C6} \\
Top-level entity & \code{Assignments $\rightarrow$ Settings $\rightarrow$ General} & the \textbf{design}, never the testbench \\
Files & \code{Project $\rightarrow$ Add/Remove Files} & design \emph{and} testbench \code{.vhd} \\
Simulator & \code{Settings $\rightarrow$ EDA Tool Settings $\rightarrow$ Simulation} & \code{Questa Altera FPGA (VHDL)} \\
Test bench & same page $\rightarrow$ \code{Compile test bench} $\rightarrow$ \code{Test Benches\ldots} $\rightarrow$ \code{New} & test bench name and top-level module = the testbench entity \\
\bottomrule
\end{tabular}
\normalsize
\end{center}

\begin{center}
\begin{tabular}{lll}
\toprule
example & top-level entity & test bench entity \\
\midrule
1 & \code{majority3} & \code{majority3\_tb} \\
2 & \code{trans\_latch} & \code{trans\_latch\_tb} \\
3 & \code{trans\_latch\_bad} & \code{sensitivity\_compare\_tb} \\
\bottomrule
\end{tabular}
\end{center}

\textbf{Why the testbench is in the file list but is not the top level.} Quartus
analyses every file in the list but only synthesises what sits underneath the
top-level entity. The testbench contains \code{wait for 10 ns}, which has no hardware
meaning; as long as it is not the top level, it is parsed and then ignored. Make it
the top level and synthesis fails.

\textbf{The flow, every time:}
\begin{enumerate}[leftmargin=2em,itemsep=1pt]
\item \code{Processing $\rightarrow$ Start $\rightarrow$ Start Analysis \& Elaboration} --- syntax, seconds
\item \code{Tools $\rightarrow$ Netlist Viewers $\rightarrow$ RTL Viewer} --- \emph{did I describe what I meant?}
\item \code{Processing $\rightarrow$ Start Compilation} (Ctrl+L) --- then \textbf{read the Messages}
\item \code{Tools $\rightarrow$ Run Simulation Tool $\rightarrow$ RTL Simulation} --- does it behave?
\end{enumerate}

\vspace{4pt}\hr

\section*{Example 1 --- Majority of three: learning the flow}

\subsection*{1.1 Why an old circuit}

You derived this circuit in Lab 0. That is deliberate: \textbf{never learn a tool and
a circuit at the same time.} If something fails today, the circuit is already known
to be right, so the problem is in how you drove the tool. You have exactly one unknown.

\subsection*{1.2 Specification, table, equation}

$Y = 1$ when at least two of $A,B,C$ are $1$.

\begin{center}
\begin{tabular}{c|ccc|c|l}
\toprule
\# & $A$ & $B$ & $C$ & $Y$ & \\
\midrule
0 & 0&0&0 & 0 & \\
1 & 0&0&1 & 0 & \\
2 & 0&1&0 & 0 & \\
3 & 0&1&1 & \textbf{1} & $m_3$ \\
4 & 1&0&0 & 0 & \\
5 & 1&0&1 & \textbf{1} & $m_5$ \\
6 & 1&1&0 & \textbf{1} & $m_6$ \\
7 & 1&1&1 & \textbf{1} & $m_7$ \\
\bottomrule
\end{tabular}
\end{center}

\begin{align*}
\text{(1)}\quad Y &= \textstyle\sum m(3,5,6,7) = \bar{A}BC + A\bar{B}C + AB\bar{C} + ABC \\
\text{(2)}\quad Y &= AB + AC + BC
  && \text{\small idempotent, distributive, complement} \\
\text{(3)}\quad & \code{Y <= (A and B) or (A and C) or (B and C);}
\end{align*}

The testbench file's header abbreviates step (2) as ``each pair absorbs $ABC$''.
In your report, name the laws as above --- the full derivation is in the Lab 0
solution, Exercise 1.

\subsection*{1.3 Anatomy of the design file}

\begin{lstlisting}
library ieee;                       -- where the standard packages live
use ieee.std_logic_1164.all;        -- gives std_logic and and/or/not

entity majority3 is                 -- the INTERFACE: what the block looks like outside
    port ( A, B, C : in  std_logic;
           Y       : out std_logic );
end majority3;

architecture behavioral of majority3 is   -- the IMPLEMENTATION: what is inside
begin
    Y <= (A and B) or (A and C) or (B and C);
end architecture behavioral;
\end{lstlisting}

\textbf{Entity = the pins. Architecture = the wiring.} One entity can have several
architectures (Lab 0 had behavioural and structural for the full adder); the entity
is the contract they all honour.

\code{<=} outside a process is a \textbf{concurrent} assignment: it is not a line
that runs once, it is a wire that is always true. Change \code{A} at any time and
\code{Y} follows.

\subsection*{1.4 Anatomy of the testbench --- the pattern for the whole semester}

\begin{lstlisting}
entity majority3_tb is              -- 1. NO PORTS: a closed world
end entity majority3_tb;

architecture test of majority3_tb is
    signal A, B, C : std_logic := '0';
    signal Y       : std_logic;

    type vector_t is record  a, b, c, y : std_logic;  end record;
    type vector_array is array (natural range <>) of vector_t;
    constant VECTORS : vector_array := (      -- 2. THE TRUTH TABLE, AS DATA
        ('0','0','0','0'), ('0','0','1','0'), ('0','1','0','0'), ('0','1','1','1'),
        ('1','0','0','0'), ('1','0','1','1'), ('1','1','0','1'), ('1','1','1','1') );
begin
    DUT : entity work.majority3                -- 3. INSTANTIATE, named association
        port map (A => A, B => B, C => C, Y => Y);

    stimulus : process
        variable errors : natural := 0;
    begin
        for i in VECTORS'range loop
            A <= VECTORS(i).a;  B <= VECTORS(i).b;  C <= VECTORS(i).c;
            wait for 10 ns;                    -- 4. LET THE CIRCUIT RESPOND
            if Y /= VECTORS(i).y then          -- 5. CHECK, don't eyeball
                report "FAIL: row " & integer'image(i) severity error;
                errors := errors + 1;
            end if;
        end loop;
        if errors = 0 then                     -- 6. VERDICT, always printed
            report "PASS - majority3: all 8 rows correct." severity note;
        end if;
        wait;                                  -- 7. STOP, or it loops forever
    end process stimulus;
end architecture test;
\end{lstlisting}

Four things in there are not obvious, and each one bites someone every semester:

\begin{itemize}[leftmargin=1.6em,itemsep=2pt]
\item \textbf{Why \code{wait for 10 ns} comes before the check.} A signal assignment
inside a process does not take effect until the process suspends. Check \code{Y} on
the line right after \code{A <= \ldots} and you are reading the \emph{previous}
vector's answer. The wait is what lets the new inputs arrive.
\item \textbf{Why the table is a \code{constant}.} Eight literal rows, copied straight
from Section 1.2 and indexed by row number, so a failure message names the
\emph{table row}. Reports say ``row 5 failed'', which you can look up; they never
say ``something went wrong around 50 ns''. By Lab 11 this table has 8 rows and 14
columns and it is still the same loop.
\item \textbf{Why the final \code{wait;}.} A process without a sensitivity list
restarts from the top when it reaches the end. Without the last \code{wait}, the
testbench reruns forever and the simulation never stops.
\item \textbf{Why \code{PASS} is printed.} A testbench that prints nothing on success
looks exactly like one that never ran.
\end{itemize}

\subsection*{1.5 Result}

\begin{lstlisting}
# ** Note: PASS - majority3: all 8 rows correct.
#    Time: 80 ns  Iteration: 0  Instance: /majority3_tb
# Errors: 0, Warnings: 0
\end{lstlisting}

\textbf{Read the time as a check:} 80~ns = 8 vectors $\times$ 10~ns. If your run ends
anywhere else, your loop did not cover the table.

\subsection*{1.6 Quartus evidence}

\begin{center}
\begin{tabular}{ll}
\toprule
Total pins & 4 \quad (3 in, 1 out) \\
Total registers & \textbf{0} \\
Logic utilisation & 1 ALM \\
Timing Analyzer & \code{Warning (332068): No clocks defined in design.} \\
& \code{Info (332140): No fmax paths to report} \\
\bottomrule
\end{tabular}
\end{center}

All three lines are \textbf{correct results}, not problems. Zero registers because
nothing is stored; one ALM because any function of up to six inputs fits in one; and
no $F_{max}$ because there is no clock to measure a frequency against. You will also
see \code{Critical Warning (332012)} saying no \code{.sdc} constraints file exists ---
true, harmless, and expected until a design has a clock.

\textbf{RTL Viewer:} three AND gates into one OR. It should look like the equation.
If it does, the tool understood you.

\vspace{4pt}\hr

\section*{Example 2 --- The transparent latch: memory}

\subsection*{2.1 Specification and table --- Type T}

\code{En} = 1: \code{Q} follows \code{D}. \quad \code{En} = 0: \code{Q} keeps its
value and ignores \code{D}.

\begin{center}
\begin{tabular}{c|cc|c|l}
\toprule
\# & \code{En} & \code{D} & $Q^{+}$ & \\
\midrule
0 & 0 & -- & $\mathbf{Q}$ & \textbf{hold} --- the memory row \\
1 & 1 & 0 & 0 & transparent \\
2 & 1 & 1 & 1 & transparent \\
\bottomrule
\end{tabular}
\end{center}

Row 0 has $Q$ in the output column. \textbf{That is the signature of memory}: the
next value depends on the present value.

\subsection*{2.2 Boolean equation}

\textbf{(1) Canonical}, taking the rows where $Q^{+}$ can be 1:
\[
Q^{+} = \underbrace{\code{En}\cdot D}_{\text{row 2}} \;+\; \underbrace{\overline{\code{En}}\cdot Q}_{\text{row 0 holding a 1}}
\]

Nothing simplifies. But read it: this is \textbf{exactly the Lab 0 multiplexer}
$\bar{S}D_0 + SD_1$, with $S = \code{En}$, $D_1 = D$ and $D_0 = Q$.

\begin{quote}
\textbf{A latch is a 2-to-1 multiplexer whose output is wired back to one of its own
inputs.} Select $D$ and new data flows in; select $Q$ and the old value circulates.
\end{quote}

Lab 0 had no output on the right-hand side of any equation. This one does, and that
single feedback path is the whole difference between logic and memory.

\textbf{(2) VHDL:}
\begin{lstlisting}
process (En, D)            -- every signal the process READS
begin
    if En = '1' then
        Q <= D;            -- rows 1, 2
    end if;                -- no else: row 0
end process;
\end{lstlisting}

\textbf{The missing \code{else} is row 0.} In a process, a signal not assigned on
some path keeps its previous value. When \code{En = '0'} nothing assigns \code{Q}, so
\code{Q} holds --- and synthesis builds storage to make that true.

\subsection*{2.3 The testbench, as a timeline}

The testbench adds one new tool: a \code{check} procedure, so the comparison and the
error message are written once and called four times.

\begin{center}
\begin{tabular}{r|cc|c|l}
\toprule
$t$ (ns) & \code{En} & \code{D} & \code{Q} checked & what it proves \\
\midrule
0  & 1 & 0 & & open, with 0 \\
10 & 1 & \textbf{1} & 0 at $t{=}10$ \ok & check 1: row 1 \\
20 & \textbf{0} & 1 & 1 at $t{=}20$ \ok & check 2: row 2 --- and \textbf{close} \\
30 & 0 & \textbf{0} & & change \code{D} \emph{while closed} \\
40 & \textbf{1} & 0 & \textbf{1} at $t{=}40$ \ok & check 3: \textbf{row 0, the hold} \\
50 & & & 0 at $t{=}50$ \ok & check 4: reopening catches up \\
\bottomrule
\end{tabular}
\end{center}

\begin{lstlisting}
# ** Note: PASS - trans_latch: all checks passed.
#    Time: 50 ns  Iteration: 0  Instance: /trans_latch_tb
# Errors: 0, Warnings: 0
\end{lstlisting}

\textbf{Two details in that timeline are design decisions:}

\textbf{Check 3 changes \code{D} after closing, not before.} Leaving \code{D}
unchanged would prove nothing --- a plain wire holds a value that nobody changes.
The hold is only demonstrated by moving the input underneath it.

\textbf{\code{En} falls at 20~ns and \code{D} changes at 30~ns --- never at the same
instant.} In simulation a simultaneous change resolves cleanly either way. In
hardware, data changing at the same moment the latch closes is a race that can store
either value. Separate them by a full step, and keep that habit when clocks arrive in
Lab 2.

\subsection*{2.4 What each check catches --- measured}

Example 2's testbench was run against three wrong designs:

\begin{center}
\begin{tabular}{lcccc|l}
\toprule
design under test & chk 1 & chk 2 & chk 3 & chk 4 & transcript \\
\midrule
correct latch & \ok & \ok & \ok & \ok & \code{PASS} \\
wire: \code{Q <= D;} & \ok & \ok & \bad & \ok & 1 failed \\
AND gate: \code{Q <= En and D;} & \ok & \ok & \bad & \ok & 1 failed \\
broken sensitivity list (Example 3) & \ok & \bad & \bad & \ok & 2 failed \\
\bottomrule
\end{tabular}
\end{center}

\textbf{Look at the wire and the AND gate:} they pass three checks out of four, and
the only one that catches them is the hold check. Neither has any memory, and every
other vector is satisfiable without memory. \textbf{Delete check 3 and a piece of
wire passes your testbench.} That is why the assignment's rubric looks for your hold
vector before anything else.

The AND gate is not hypothetical: it is exactly what you get by adding
\code{else Q <= '0';} to the process. It is complete, warning-free, and not a latch.

\subsection*{2.5 Quartus evidence}

\textbf{Messages.} One warning, and it is the one you want:
\begin{lstlisting}
Warning (10631): VHDL Process Statement warning at trans_latch.vhd(35):
  inferring latch(es) for signal or variable "Q", which holds its previous
  value in one or more paths through the process
\end{lstlisting}

This is the ``inferred latch'' warning. In almost every other design you write it
means you forgot a branch. \textbf{Here it confirms the design is right} --- and
having caused it on purpose, you will recognise it instantly when it is an accident
in Lab 3.

\textbf{Inferred-latch table} (Compilation Report $\rightarrow$ Analysis \& Synthesis):
\begin{center}
\begin{tabular}{lll}
\toprule
Latch name & Latch enable signal & Free of timing hazards \\
\midrule
\code{Q\$latch} & \code{En} & yes \\
\bottomrule
\end{tabular}
\end{center}

Quartus worked out on its own that \code{En} controls the storage. It agrees with
your table.

\textbf{Flow summary and netlist.}
\begin{center}
\begin{tabular}{ll}
\toprule
Total pins & 3 \\
Total registers & \textbf{0} \quad --- a latch is not a register \\
Post-synthesis netlist & 1 LUT with \textbf{3 data inputs} \\
Timing Analyzer & clock \code{En} derived; min.\ pulse width slack $+0.365$~ns \\
\bottomrule
\end{tabular}
\end{center}

\textbf{The Timing Analyzer invents a clock --- read this before you panic.} With no
\code{.sdc} constraints file, Quartus looks for anything that controls storage and
assumes it is a clock with a 1.000~ns period:
\begin{lstlisting}
Info (332105): create_clock -period 1.000 -name En En
Info (332140): No fmax paths to report
\end{lstlisting}
That is a 1~GHz requirement nobody asked for. There are \textbf{no setup or hold
paths} (nothing goes from one storage element to another), so there is \textbf{no
$F_{max}$}. The one check left is \emph{minimum pulse width}: can \code{En} stay high
long enough at 1~GHz? Here it can, by 0.365~ns.

On another build it may not. The 8-bit reference \code{TransLatch8} misses it by
38~ps at the fast corner and prints \code{Critical Warning (332148): Timing
requirements not met}. \textbf{Against an invented 1~ns clock, that warning means
nothing.} Report it honestly and say why; a meaningful timing number needs a real
clock and a real constraint, and both arrive in Lab 2.

\textbf{Three data inputs for a circuit with two inputs.} The third LUT input is
\code{Q} itself, fed back. You are reading the feedback term $\overline{\code{En}}\cdot Q$
straight out of the netlist statistics.

\textbf{RTL Viewer:} a latch symbol with \code{D} on its data pin and \code{En} on
its enable. No loop of gates --- Quartus draws the primitive, not the multiplexer.

\vspace{4pt}\hr

\section*{Example 3 --- One missing name}

\subsection*{3.1 The only difference}

\begin{center}
\begin{tabular}{ll}
\toprule
\code{trans\_latch} (Example 2) & \code{trans\_latch\_bad} \\
\midrule
\code{process (En, D)} & \code{process (En)} \\
\bottomrule
\end{tabular}
\end{center}

The process body is identical. \textbf{Rule: the sensitivity list must name every
signal the process reads.} This process reads \code{En} and \code{D}.

\subsection*{3.2 Simulation --- the side-by-side testbench}

Both latches receive the same stimulus:

\begin{center}
\begin{tabular}{r|cc|cc|l}
\toprule
$t$ (ns) & \code{En} & \code{D} & \code{Q\_good} & \code{Q\_bad} & \\
\midrule
0  & \textbf{1} & 0 & 0 & 0 & \code{En} changed: \textbf{both} processes wake \\
10 & 1 & \textbf{1} & \textbf{1} & \textbf{0} & only \code{D} changed: \textbf{only the good process wakes} \\
20 & 1 & 1 & 1 & 0 & assertions evaluated \\
\bottomrule
\end{tabular}
\end{center}

\begin{lstlisting}
# ** Error: EXPECTED FAILURE: the broken latch did not follow D.
#           One missing name in a sensitivity list caused this.
#    Time: 20 ns  Iteration: 0  Instance: /sensitivity_compare_tb
# Errors: 1, Warnings: 0
\end{lstlisting}

\textbf{Exactly one error, and it is the intended one.} If you see two, the
\emph{good} latch is also broken --- go back to Example 2.

\textbf{Waveform:} look between 10 and 20~ns. \code{D} rises, \code{Q\_good} rises
with it, \code{Q\_bad} stays low. (The comment at the top of
\code{sensitivity\_compare\_tb.vhd} says \code{D} changes at 20~ns; it changes at
10~ns and is checked at 20~ns.)

\subsection*{3.3 What the broken design actually is}

Trace the broken process. It runs only when \code{En} changes. When it runs, it
copies \code{D} only if \code{En} is now \code{'1'}. So \code{Q} is updated
\textbf{only at the instant \code{En} goes from 0 to 1}, and holds at every other
time.

\begin{quote}
\textbf{That is not a broken latch. It is a flip-flop clocked by the rising edge of
\code{En}.} One missing name changed which kind of storage you described.
\end{quote}

You build a flip-flop on purpose in Lab 2 --- with \code{rising\_edge(clk)}, which
says so explicitly. Example 3 builds one by accident.

\subsection*{3.4 Synthesis --- the part most textbooks get wrong}

The usual story is that synthesis ignores the sensitivity list, builds the correct
latch, and prints a warning. \textbf{For this exact code, Quartus 25.1 does not do
that.} Measured:

\begin{center}
\begin{tabular}{lcc}
\toprule
& \code{trans\_latch} & \code{trans\_latch\_bad} \\
\midrule
Warnings & \code{10631} (latch inferred) & \textbf{none} \\
Inferred latches & 1 & 0 \\
Total registers & 0 & \textbf{1} \\
Netlist & 1 LUT & \textbf{1 flip-flop} (\code{arriav\_ff}) \\
Min.\ pulse width on \code{En} & $+0.365$~ns & $\mathbf{-0.394}$~ns, \code{332148} \\
\bottomrule
\end{tabular}
\end{center}

Both designs get the same invented 1~ns clock on \code{En}. The flip-flop fails the
pulse-width check and the latch passes it --- a faint but real sign in the timing
report that the two storage elements are different primitives.

A process that wakes only on \code{En} and then tests \code{En = '1'} matches an old
flip-flop template that synthesis tools recognise, so Quartus reads the sensitivity
list as part of the clock description and builds a \textbf{flip-flop, silently}.
Simulation and hardware \emph{agree} --- and both are wrong.

\textbf{Now compare your assignment's starter}, which has the same missing names but
a \code{Clr} branch in front. The \code{Clr} test breaks the flip-flop template, so
Quartus falls back to ignoring the list:

\begin{center}
\small
\begin{tabular}{>{\raggedright\arraybackslash}p{6.2cm}>{\raggedright\arraybackslash}p{4.6cm}>{\raggedright\arraybackslash}p{4.6cm}}
\toprule
code shape & Quartus says & result \\
\midrule
\code{process (En)}, \code{if En='1'} \newline (Example 3)
  & nothing
  & flip-flop in sim \emph{and} hardware \\
\code{process (En)}, \code{if Clr \ldots elsif En} \newline (your starter)
  & \code{10492} (\code{Clr}), \code{10492} (\code{D}), \code{10631}
  & latch in hardware, broken simulation \\
\code{process (En, D)}, \code{if Clr \ldots elsif En} \newline (check-off question 2)
  & \code{10492} (\code{Clr}), \code{10631}
  & latch in hardware, simulation misses the clear \\
\code{process (Clr, En, D)} \newline (correct)
  & \code{10631} only
  & latch in both \\
\bottomrule
\end{tabular}
\normalsize
\end{center}

Warning 10492 reads:
\begin{lstlisting}
Warning (10492): VHDL Process Statement warning at TransLatch.vhd(33): signal "Clr"
  is read inside the Process Statement but isn't in the Process Statement's
  sensitivity list
\end{lstlisting}

\textbf{The same bug produces two different failures depending on the shape of the
code around it.} In rows 2 and 3 you tested a different circuit from the one you
built --- the simulation/synthesis mismatch. In row 1 there is no mismatch at all,
and no warning, and the design is still wrong.

\subsection*{3.5 What caught it --- and why the course asks for all the evidence}

\begin{center}
\begin{tabular}{lccc}
\toprule
evidence source & Example 3 (flip-flop) & starter (mismatch) \\
\midrule
self-checking testbench & \textbf{caught} (1 error) & \textbf{caught} \\
Messages / warnings & silent & \textbf{caught} (10492) \\
Flow Summary: registers & \textbf{caught} (1, expected 0) & silent (0) \\
RTL Viewer & \textbf{caught} (flip-flop symbol, \code{En} on clock) & silent (latch) \\
\bottomrule
\end{tabular}
\end{center}

No single source catches both. The testbench catches both only because it contains a
vector that changes \code{D} while \code{En} is already high. \textbf{That is the
argument for reading every piece of evidence every time}: each one is blind to
something another one sees.

\textbf{Two habits from this example:}
\begin{enumerate}[leftmargin=2em,itemsep=1pt]
\item \textbf{Silence is not correctness.} The correct latch prints a warning. The
broken one prints nothing.
\item \textbf{Know the number you expect before you open the report.} ``Registers: 0''
is a prediction you make from your table. When the report says 1, you have found a
bug without running a single vector.
\end{enumerate}

\vspace{4pt}\hr

\section*{Summary --- the three examples in Quartus}

\begin{center}
\begin{tabular}{lccc}
\toprule
& \code{majority3} & \code{trans\_latch} & \code{trans\_latch\_bad} \\
\midrule
pins & 4 & 3 & 3 \\
registers & 0 & 0 & \textbf{1} \\
inferred latches & 0 & 1 & 0 \\
warnings (design) & none & 10631 & none \\
Timing Analyzer & no clocks & \code{En} derived, passes & \code{En} derived, \textbf{fails} \\
simulation & PASS, 80 ns & PASS, 50 ns & 1 expected error \\
what it is & logic & latch & \textbf{flip-flop on \code{En}} \\
\bottomrule
\end{tabular}
\end{center}

\textbf{Into the assignment:} \code{TransLatch} is Example 2 with an asynchronous
clear added in front --- which is precisely what changes Example 3's silent flip-flop
into the warned, mismatched latch your starter will give you. Fix the sensitivity
list to \code{(Clr, En, D)}, keep the missing \code{else}, and keep a hold vector
that moves \code{D} underneath a closed latch.

\end{document}
